\documentclass[11pt]{article}
\usepackage[utf8]{inputenc}
\usepackage{amsmath,amssymb,amsfonts,amsthm}
\usepackage{geometry}
\geometry{top=1in,bottom=1in,left=1in,right=1in}
\usepackage{xcolor}
\usepackage{enumitem}
\usepackage{tcolorbox}
\tcbuselibrary{skins,breakable}
\usepackage{microtype}
\usepackage{hyperref}

% --- Color Palette (sleek mathematical theme) ---
\definecolor{primary}{HTML}{1A365D}      % Deep Blue
\definecolor{secondary}{HTML}{2B6CB0}    % Slate Blue
\definecolor{accent}{HTML}{C53030}       % Muted Red
\definecolor{lightbg}{HTML}{F7FAFC}      % Off-white
\definecolor{boxborder}{HTML}{E2E8F0}    % Light Gray

\hypersetup{
    colorlinks=true,
    linkcolor=secondary,
    filecolor=accent,      
    urlcolor=secondary,
}

% --- Custom Environments ---
\newtcolorbox[auto counter,number within=section]{examplebox}[2][]{
    colback=lightbg,
    colframe=secondary,
    arc=1mm,
    outer arc=1mm,
    left=3mm,
    right=3mm,
    top=3mm,
    bottom=3mm,
    boxrule=0.5mm,
    fonttitle=\bfseries,
    coltitle=white,
    colbacktitle=secondary,
    title={Example \thetcbcounter: #2},
    breakable,
    #1
}

\newtcolorbox[auto counter,number within=section]{formulabox}[1][]{
    colback=white,
    colframe=primary,
    arc=1mm,
    outer arc=1mm,
    left=3mm,
    right=3mm,
    top=3mm,
    bottom=3mm,
    boxrule=0.7mm,
    fonttitle=\bfseries,
    coltitle=white,
    colbacktitle=primary,
    title={Key Formula(s)},
    breakable,
    #1
}

\newtcolorbox[auto counter,number within=section]{theorembox}[1][]{
    colback=lightbg,
    colframe=accent,
    arc=1mm,
    outer arc=1mm,
    left=3mm,
    right=3mm,
    top=3mm,
    bottom=3mm,
    boxrule=0.5mm,
    fonttitle=\bfseries,
    coltitle=white,
    colbacktitle=accent,
    title={Theorem},
    breakable,
    #1
}

\title{\Huge\bfseries\color{primary} MTH253 -- Calculus III Notes \\ \Large Chapters 1 \& 2: Multivariable Differentiation \& Multiple Integration}
\author{\large VNU-HCM University of Science}
\date{\large \today}

\begin{document}

\maketitle

\tableofcontents
\newpage

% =========================================================================
% CHAPTER 1
% =========================================================================
\section{Chapter 1: Multivariable Differentiation}

\subsection{Coordinates and Surfaces in \texorpdfstring{$\mathbb{R}^3$}{R3}}
\subsubsection*{Key Notations}
\begin{itemize}
    \item $\mathbb{R}^3$: Three-dimensional real coordinate space, representing points $(x,y,z)$.
    \item $d(P_1, P_2)$: The Euclidean distance between points $P_1$ and $P_2$.
    \item $S$: A surface represented by a single equation in $x, y, z$.
\end{itemize}

\subsubsection*{Concise Explanation \& Formulas}
In 3D space, an equation with one constraint generally defines a surface.
\begin{formulabox}
    \textbf{Distance Formula in $\mathbb{R}^3$:}
    \[ d(P_1, P_2) = \sqrt{(x_2-x_1)^2 + (y_2-y_1)^2 + (z_2-z_1)^2} \]
    \textbf{Equation of a Sphere:} A sphere with center $C(h,k,l)$ and radius $R$:
    \[ (x-h)^2 + (y-k)^2 + (z-l)^2 = R^2 \]
    \textbf{Cylindrical Surfaces:} Surfaces whose equations contain only two variables (e.g., $x^2 + y^2 = R^2$) project as a curve in 2D and extend infinitely parallel to the axis of the missing variable ($z$-axis here) in $\mathbb{R}^3$.
\end{formulabox}

\begin{examplebox}{Distance and Cylindrical Surfaces}
    \textbf{Part A: Distance and Sphere}\\
    Find the equation of the sphere centered at $C(2, -1, 3)$ that passes through $P(4, 2, -1)$.\\
    \textit{Solution:}\\
    First, find the radius $R$ as the distance from $C$ to $P$:
    \[ R = d(C,P) = \sqrt{(4-2)^2 + (2-(-1))^2 + (-1-3)^2} = \sqrt{2^2 + 3^2 + (-4)^2} = \sqrt{4 + 9 + 16} = \sqrt{29} \]
    The equation of the sphere is:
    \[ (x-2)^2 + (y+1)^2 + (z-3)^2 = 29 \]
    
    \textbf{Part B: Cylindrical Surface}\\
    Describe the surface represented by $x^2 + y^2 = 1$ in $\mathbb{R}^3$.\\
    \textit{Solution:}\\
    The equation $x^2 + y^2 = 1$ does not restrict $z$. In the $xy$-plane, this is a circle of radius 1 centered at the origin. Since $z$ is free, we replicate this circle at all values of $z$, which forms a circular cylinder of radius 1 whose central axis is the $z$-axis.
\end{examplebox}


\subsection{Functions of Several Variables}
\subsubsection*{Key Notations}
\begin{itemize}
    \item $f: D \to \mathbb{R}$: A function mapping a domain $D \subset \mathbb{R}^n$ to $\mathbb{R}$.
    \item $D(f)$: The domain of $f$, the set of all inputs for which $f$ is well-defined.
    \item Level Curves: Contours in 2D where $f(x,y) = k$ (for constant $k$).
    \item Level Surfaces: Surfaces in 3D where $f(x,y,z) = k$.
\end{itemize}

\subsubsection*{Concise Explanation \& Formulas}
A multivariable function maps a tuple of variables to a single real output. The graph of $z = f(x,y)$ is a surface in $\mathbb{R}^3$. 
\begin{formulabox}
    \textbf{Level Curves:} The set of points in the domain of $f(x,y)$ satisfying:
    \[ \{(x,y) \in D \mid f(x,y) = k\}, \quad k \in \text{Range}(f) \]
    \textbf{Level Surfaces:} The set of points in the domain of $f(x,y,z)$ satisfying:
    \[ \{(x,y,z) \in D \mid f(x,y,z) = k\} \]
\end{formulabox}

\begin{examplebox}{Domain and Level Curves}
    Let $f(x,y) = \sqrt{9 - x^2 - y^2}$.\\
    1. Find the domain and range of $f$.\\
    2. Sketch the level curves for $k = 0, \sqrt{5}, 3$.\\
    \textit{Solution:}\\
    1. For the square root to be defined, we must have $9 - x^2 - y^2 \ge 0 \implies x^2 + y^2 \le 9$. Thus, the domain $D$ is the closed disk centered at the origin with radius 3:
    \[ D = \{(x,y) \in \mathbb{R}^2 \mid x^2 + y^2 \le 9\} \]
    The range is $[0, 3]$ because $0 \le \sqrt{9 - x^2 - y^2} \le \sqrt{9} = 3$.\\
    2. The level curves are given by $\sqrt{9 - x^2 - y^2} = k \implies x^2 + y^2 = 9 - k^2$.
    \begin{itemize}
        \item For $k=0$: $x^2 + y^2 = 9$ (circle of radius 3).
        \item For $k=\sqrt{5}$: $x^2 + y^2 = 9 - 5 = 4$ (circle of radius 2).
        \item For $k=3$: $x^2 + y^2 = 9 - 9 = 0$ (the origin $(0,0)$).
    \end{itemize}
\end{examplebox}


\subsection{Limits and Continuity}
\subsubsection*{Key Notations}
\begin{itemize}
    \item $\lim\limits_{(x,y) \to (a,b)} f(x,y) = L$: The limit of $f$ as $(x,y)$ approaches $(a,b)$ is $L$.
    \item $C_1, C_2$: Paths of approach in the input plane.
\end{itemize}

\subsubsection*{Concise Explanation \& Formulas}
For a limit to exist in $\mathbb{R}^2$, the function must approach the same value $L$ along \textit{every possible path} approaching $(a,b)$. If we find two paths yielding different limits, the limit does not exist.
\begin{formulabox}
    \textbf{Limit Non-Existence Test:} If $f(x,y) \to L_1$ along path $C_1$ and $f(x,y) \to L_2$ along path $C_2$ as $(x,y) \to (a,b)$, and $L_1 \neq L_2$, then:
    \[ \lim_{(x,y) \to (a,b)} f(x,y) \quad \text{does not exist.} \]
\end{formulabox}

\begin{theorembox}[title={Squeeze Theorem}]
    If $g(x,y) \le f(x,y) \le h(x,y)$ for all $(x,y)$ near $(a,b)$ and $\lim\limits_{(x,y) \to (a,b)} g(x,y) = \lim\limits_{(x,y) \to (a,b)} h(x,y) = L$, then:
    \[ \lim_{(x,y) \to (a,b)} f(x,y) = L \]
\end{theorembox}

\begin{examplebox}{Evaluating Limits and Non-Existence}
    \textbf{Part A: Non-existence}\\
    Show that $\lim\limits_{(x,y) \to (0,0)} \frac{x^2 - y^2}{x^2 + y^2}$ does not exist.\\
    \textit{Solution:}\\
    Test path $C_1$ along the $x$-axis ($y = 0$):
    \[ \lim_{(x,0) \to (0,0)} \frac{x^2 - 0^2}{x^2 + 0^2} = \lim_{x \to 0} \frac{x^2}{x^2} = 1 \]
    Test path $C_2$ along the $y$-axis ($x = 0$):
    \[ \lim_{(0,y) \to (0,0)} \frac{0^2 - y^2}{0^2 + y^2} = \lim_{y \to 0} \frac{-y^2}{y^2} = -1 \]
    Since the two limits are different ($1 \neq -1$), the limit does not exist.
    
    \textbf{Part B: Squeeze Theorem}\\
    Evaluate $\lim\limits_{(x,y) \to (0,0)} \frac{3x^2y}{x^2 + y^2}$.\\
    \textit{Solution:}\\
    Notice that $x^2 \le x^2 + y^2 \implies \frac{x^2}{x^2+y^2} \le 1$. Taking absolute values:
    \[ 0 \le \left|\frac{3x^2y}{x^2 + y^2}\right| = 3|y|\frac{x^2}{x^2+y^2} \le 3|y| \]
    Therefore:
    \[ -3|y| \le \frac{3x^2y}{x^2 + y^2} \le 3|y| \]
    Since $\lim\limits_{(x,y)\to(0,0)} (-3|y|) = 0$ and $\lim\limits_{(x,y)\to(0,0)} 3|y| = 0$, by the Squeeze Theorem:
    \[ \lim_{(x,y) \to (0,0)} \frac{3x^2y}{x^2 + y^2} = 0 \]
\end{examplebox}


\subsection{Partial Derivatives}
\subsubsection*{Key Notations}
\begin{itemize}
    \item $f_x = \frac{\partial f}{\partial x}$: Partial derivative of $f$ with respect to $x$ (holding $y$ constant).
    \item $f_y = \frac{\partial f}{\partial y}$: Partial derivative of $f$ with respect to $y$ (holding $x$ constant).
    \item $f_{xx}, f_{yy}, f_{xy}, f_{yx}$: Second-order partial derivatives.
\end{itemize}

\subsubsection*{Concise Explanation \& Formulas}
Partial derivatives measure the rates of change of a multivariable function along the coordinate directions.
\begin{formulabox}
    \textbf{First-Order Partial Derivatives:}
    \[ f_x(x,y) = \lim_{h \to 0} \frac{f(x+h,y) - f(x,y)}{h}, \quad f_y(x,y) = \lim_{h \to 0} \frac{f(x,y+h) - f(x,y)}{h} \]
    \textbf{Second-Order Mixed Partials:}
    \[ f_{xy} = \frac{\partial^2 f}{\partial y \partial x} = \frac{\partial}{\partial y}\left(\frac{\partial f}{\partial x}\right), \quad f_{yx} = \frac{\partial^2 f}{\partial x \partial y} = \frac{\partial}{\partial x}\left(\frac{\partial f}{\partial y}\right) \]
\end{formulabox}

\begin{theorembox}[title={Clairaut's Theorem}]
    If $f$ is defined on a disk $D$ containing $(a,b)$, and mixed partial derivatives $f_{xy}$ and $f_{yx}$ are both continuous on $D$, then:
    \[ f_{xy}(a,b) = f_{yx}(a,b) \]
\end{theorembox}

\begin{examplebox}{First and Second Order Partials}
    Let $f(x,y) = x^3 + x^2y^3 - 2y^2$. Find all first and second order partial derivatives, and verify Clairaut's Theorem.\\
    \textit{Solution:}\\
    First-order partial derivatives:
    \[ f_x = \frac{\partial}{\partial x}(x^3 + x^2y^3 - 2y^2) = 3x^2 + 2xy^3 \]
    \[ f_y = \frac{\partial}{\partial y}(x^3 + x^2y^3 - 2y^2) = 3x^2y^2 - 4y \]
    Second-order partial derivatives:
    \[ f_{xx} = \frac{\partial}{\partial x}(3x^2 + 2xy^3) = 6x + 2y^3 \]
    \[ f_{yy} = \frac{\partial}{\partial y}(3x^2y^2 - 4y) = 6x^2y - 4 \]
    Mixed partial derivatives:
    \[ f_{xy} = \frac{\partial}{\partial y}(3x^2 + 2xy^3) = 6xy^2 \]
    \[ f_{yx} = \frac{\partial}{\partial x}(3x^2y^2 - 4y) = 6xy^2 \]
    Since $f_{xy} = f_{yx} = 6xy^2$ and this expression is continuous everywhere, Clairaut's Theorem is verified.
\end{examplebox}


\subsection{Tangent Planes and Linearization}
\subsubsection*{Key Notations}
\begin{itemize}
    \item $L(x,y)$: Linearization of $f(x,y)$ at $(a,b)$.
    \item $dz$: The total differential of $z = f(x,y)$.
\end{itemize}

\subsubsection*{Concise Explanation \& Formulas}
Near a point where a function is differentiable, the surface $z=f(x,y)$ can be approximated by its tangent plane. The equation of the tangent plane defines the linear approximation (linearization) of the function.
\begin{formulabox}
    \textbf{Tangent Plane Equation:} The tangent plane to $z = f(x,y)$ at $P(x_0, y_0, z_0)$ is:
    \[ z - z_0 = f_x(x_0, y_0)(x - x_0) + f_y(x_0, y_0)(y - y_0) \]
    \textbf{Linearization and Linear Approximation:}
    \[ L(x,y) = f(a,b) + f_x(a,b)(x-a) + f_y(a,b)(y-b) \implies f(x,y) \approx L(x,y) \]
    \textbf{Total Differential:} For $z = f(x,y)$, the differential $dz$ is:
    \[ dz = \frac{\partial z}{\partial x}dx + \frac{\partial z}{\partial y}dy = f_x(x,y)dx + f_y(x,y)dy \]
\end{formulabox}

\begin{examplebox}{Tangent Plane and Linearization}
    For $f(x,y) = 2x^2 + y^2$, find the tangent plane and the linearization at $(1, 1)$. Then approximate $f(1.05, 0.95)$.\\
    \textit{Solution:}\\
    First, find the function value and partial derivatives at $(1,1)$:
    \[ f(1,1) = 2(1)^2 + (1)^2 = 3 \implies z_0 = 3 \]
    \[ f_x = 4x \implies f_x(1,1) = 4 \]
    \[ f_y = 2y \implies f_y(1,1) = 2 \]
    The tangent plane equation is:
    \[ z - 3 = 4(x - 1) + 2(y - 1) \implies z = 4x + 2y - 3 \]
    The linearization is:
    \[ L(x,y) = 3 + 4(x - 1) + 2(y - 1) \]
    We approximate $f(1.05, 0.95)$ using $L(1.05, 0.95)$:
    \[ f(1.05, 0.95) \approx 3 + 4(1.05 - 1) + 2(0.95 - 1) = 3 + 4(0.05) + 2(-0.05) = 3 + 0.20 - 0.10 = 3.10 \]
    (The exact value is $2(1.05)^2 + (0.95)^2 = 2.205 + 0.9025 = 3.1075$, showing the approximation is very close.)
\end{examplebox}


\subsection{The Chain Rule and Implicit Differentiation}
\subsubsection*{Key Notations}
\begin{itemize}
    \item Tree Diagram: Used to trace dependency of variables.
    \item $\frac{\partial z}{\partial s}$, $\frac{\partial z}{\partial t}$: Partial derivatives when variables depend on multiple parameters.
\end{itemize}

\subsubsection*{Concise Explanation \& Formulas}
When variables themselves are functions of other variables, we use the chain rule. We sum the products of derivatives along all paths in the dependency tree.
\begin{formulabox}
    \textbf{Chain Rule: Case 1 (One Parameter):} If $z = f(x,y)$ and $x=x(t), y=y(t)$, then:
    \[ \frac{dz}{dt} = \frac{\partial z}{\partial x}\frac{dx}{dt} + \frac{\partial z}{\partial y}\frac{dy}{dt} \]
    \textbf{Chain Rule: Case 2 (Two Parameters):} If $z = f(x,y)$ and $x=x(s,t), y=y(s,t)$, then:
    \[ \frac{\partial z}{\partial s} = \frac{\partial z}{\partial x}\frac{\partial x}{\partial s} + \frac{\partial z}{\partial y}\frac{\partial y}{\partial s}, \quad \frac{\partial z}{\partial t} = \frac{\partial z}{\partial x}\frac{\partial x}{\partial t} + \frac{\partial z}{\partial y}\frac{\partial y}{\partial t} \]
    \textbf{Implicit Differentiation:} If $F(x,y,z) = 0$ defines $z = f(x,y)$ implicitly, then:
    \[ \frac{\partial z}{\partial x} = -\frac{F_x}{F_z}, \quad \frac{\partial z}{\partial y} = -\frac{F_y}{F_z} \quad (\text{for } F_z \neq 0) \]
\end{formulabox}

\begin{examplebox}{Chain Rule and Implicit Differentiation}
    \textbf{Part A: Chain Rule}\\
    If $z = x^2y + 3xy^4$, where $x = \sin(2t)$ and $y = \cos(t)$, find $\frac{dz}{dt}$ at $t=0$.\\
    \textit{Solution:}\\
    At $t = 0$: $x = \sin(0) = 0$ and $y = \cos(0) = 1$.\\
    Find the partial derivatives of $z$:
    \[ \frac{\partial z}{\partial x} = 2xy + 3y^4 \implies \left.\frac{\partial z}{\partial x}\right|_{(0,1)} = 2(0)(1) + 3(1)^4 = 3 \]
    \[ \frac{\partial z}{\partial y} = x^2 + 12xy^3 \implies \left.\frac{\partial z}{\partial y}\right|_{(0,1)} = 0^2 + 12(0)(1)^3 = 0 \]
    Find the derivatives of $x$ and $y$ at $t = 0$:
    \[ \frac{dx}{dt} = 2\cos(2t) \implies \left.\frac{dx}{dt}\right|_{t=0} = 2, \quad \frac{dy}{dt} = -\sin(t) \implies \left.\frac{dy}{dt}\right|_{t=0} = 0 \]
    Apply Case 1 of the Chain Rule:
    \[ \frac{dz}{dt} = \frac{\partial z}{\partial x}\frac{dx}{dt} + \frac{\partial z}{\partial y}\frac{dy}{dt} = (3)(2) + (0)(0) = 6 \]
    
    \textbf{Part B: Implicit Differentiation}\\
    Find $\frac{\partial z}{\partial x}$ if $x^3 + y^3 + z^3 + 6xyz = 1$.\\
    \textit{Solution:}\\
    Define $F(x,y,z) = x^3 + y^3 + z^3 + 6xyz - 1 = 0$.\\
    Find the partial derivatives of $F$:
    \[ F_x = 3x^2 + 6yz, \quad F_z = 3z^2 + 6xy \]
    Use the implicit differentiation formula:
    \[ \frac{\partial z}{\partial x} = -\frac{F_x}{F_z} = -\frac{3x^2 + 6yz}{3z^2 + 6xy} = -\frac{x^2 + 2yz}{z^2 + 2xy} \]
\end{examplebox}


\subsection{Gradients and Directional Derivatives}
\subsubsection*{Key Notations}
\begin{itemize}
    \item $\nabla f$: The gradient vector of $f$.
    \item $D_{\vec{u}}f$: The directional derivative of $f$ in the direction of the unit vector $\vec{u}$.
\end{itemize}

\subsubsection*{Concise Explanation \& Formulas}
The gradient $\nabla f$ points in the direction of maximum rate of increase of the function. The directional derivative calculates the rate of change of $f$ along any arbitrary direction vector $\vec{u}$.
\begin{formulabox}
    \textbf{Gradient Vector:}
    \[ \nabla f(x,y) = \langle f_x, f_y \rangle = f_x\mathbf{i} + f_y\mathbf{j}, \quad \nabla f(x,y,z) = \langle f_x, f_y, f_z \rangle \]
    \textbf{Directional Derivative (using unit vector $\vec{u}$):}
    \[ D_{\vec{u}}f(x,y) = \nabla f(x,y) \cdot \vec{u} \]
    If given a general direction vector $\vec{v}$, you \textbf{must} normalize it first: $\vec{u} = \frac{\vec{v}}{\|\vec{v}\|}$.\\
    \textbf{Properties of the Gradient:}
    \begin{itemize}
        \item The maximum value of $D_{\vec{u}}f(x,y)$ is $\|\nabla f(x,y)\|$, occurring when $\vec{u}$ is in the direction of $\nabla f$.
        \item The minimum value is $-\|\nabla f(x,y)\|$, in the direction of $-\nabla f$.
        \item The gradient vector $\nabla f(x_0,y_0)$ is always orthogonal to the level curve $f(x,y)=k$ passing through $(x_0,y_0)$.
    \end{itemize}
\end{formulabox}

\begin{examplebox}{Directional Derivative and Maximum Rate of Change}
    Let $f(x,y) = x^2y^3 - 4y$.\\
    1. Find the directional derivative of $f$ at $P(2,-1)$ in the direction of $\vec{v} = \langle 2,5 \rangle$.\\
    2. Find the maximum rate of change of $f$ at $P(2,-1)$ and the direction in which it occurs.\\
    \textit{Solution:}\\
    1. Calculate the gradient:
    \[ \nabla f(x,y) = \langle 2xy^3, 3x^2y^2 - 4 \rangle \]
    At $P(2,-1)$:
    \[ \nabla f(2,-1) = \langle 2(2)(-1)^3, 3(2)^2(-1)^2 - 4 \rangle = \langle -4, 12 - 4 \rangle = \langle -4, 8 \rangle \]
    Normalize the direction vector $\vec{v}$:
    \[ \|\vec{v}\| = \sqrt{2^2 + 5^2} = \sqrt{29} \implies \vec{u} = \left\langle \frac{2}{\sqrt{29}}, \frac{5}{\sqrt{29}} \right\rangle \]
    The directional derivative is:
    \[ D_{\vec{u}}f(2,-1) = \nabla f(2,-1) \cdot \vec{u} = -4\left(\frac{2}{\sqrt{29}}\right) + 8\left(\frac{5}{\sqrt{29}}\right) = \frac{-8 + 40}{\sqrt{29}} = \frac{32}{\sqrt{29}} \]
    2. The maximum rate of change is:
    \[ \|\nabla f(2,-1)\| = \|\langle -4, 8 \rangle\| = \sqrt{(-4)^2 + 8^2} = \sqrt{16 + 64} = \sqrt{80} = 4\sqrt{5} \]
    This maximum rate of change occurs in the direction of the gradient vector $\nabla f(2,-1) = \langle -4, 8 \rangle$.
\end{examplebox}


\subsection{Extrema of Multivariable Functions}
\subsubsection*{Key Notations}
\begin{itemize}
    \item Critical Point: A point $(a,b)$ where $f_x(a,b)=0$ and $f_y(a,b)=0$ (or either is undefined).
    \item $D$: The discriminant (Hessian determinant) used in the Second Derivatives Test.
    \item Saddle Point: A critical point that is neither a local max nor local min.
\end{itemize}

\subsubsection*{Concise Explanation \& Formulas}
To locate local extrema, we find critical points and classify them using second-order partial derivatives.
\begin{formulabox}
    \textbf{Second Derivatives Test:} Let $(a,b)$ be a critical point of $f$ where second partials are continuous. Let:
    \[ D = f_{xx}(a,b)f_{yy}(a,b) - [f_{xy}(a,b)]^2 \]
    \begin{itemize}
        \item If $D > 0$ and $f_{xx}(a,b) > 0 \implies f(a,b)$ is a \textbf{local minimum}.
        \item If $D > 0$ and $f_{xx}(a,b) < 0 \implies f(a,b)$ is a \textbf{local maximum}.
        \item If $D < 0 \implies (a,b)$ is a \textbf{saddle point}.
        \item If $D = 0 \implies$ the test is \textbf{inconclusive} (requires alternative analysis).
    \end{itemize}
\end{formulabox}

\begin{examplebox}{Finding Local Extrema and Saddle Points}
    Find the local maximum and minimum values and saddle points of $f(x,y) = x^4 + y^4 - 4xy + 1$.\\
    \textit{Solution:}\\
    Step 1: Find critical points.
    \[ f_x = 4x^3 - 4y = 0 \implies y = x^3 \]
    \[ f_y = 4y^3 - 4x = 0 \implies x = y^3 \]
    Substitute $y = x^3$ into the second equation:
    \[ x = (x^3)^3 = x^9 \implies x(x^8 - 1) = 0 \]
    This gives real solutions $x = -1, 0, 1$.
    \begin{itemize}
        \item If $x=0$, $y=0 \implies (0,0)$.
        \item If $x=1$, $y=1 \implies (1,1)$.
        \item If $x=-1$, $y=-1 \implies (-1,-1)$.
    \end{itemize}
    Step 2: Apply the Second Derivatives Test.
    Calculate second partials:
    \[ f_{xx} = 12x^2, \quad f_{yy} = 12y^2, \quad f_{xy} = -4 \]
    The discriminant is $D(x,y) = f_{xx}f_{yy} - (f_{xy})^2 = (12x^2)(12y^2) - (-4)^2 = 144x^2y^2 - 16$.
    \begin{itemize}
        \item At $(0,0)$: $D(0,0) = -16 < 0$. Thus, \textbf{$(0,0)$ is a saddle point}.
        \item At $(1,1)$: $D(1,1) = 144(1)(1) - 16 = 128 > 0$ and $f_{xx}(1,1) = 12 > 0$. Thus, \textbf{$f(1,1) = -1$ is a local minimum}.
        \item At $(-1,-1)$: $D(-1,-1) = 128 > 0$ and $f_{xx}(-1,-1) = 12 > 0$. Thus, \textbf{$f(-1,-1) = -1$ is a local minimum}.
    \end{itemize}
\end{examplebox}


\subsection{Lagrange Multipliers}
\subsubsection*{Key Notations}
\begin{itemize}
    \item $f(x,y,z)$: The objective function to maximize or minimize.
    \item $g(x,y,z) = k$: The constraint equation.
    \item $\lambda$: The Lagrange multiplier (scalar).
\end{itemize}

\subsubsection*{Concise Explanation \& Formulas}
Lagrange multipliers find the extrema of an objective function subject to constraints. Geometrically, at an extreme point, the level curves of $f$ and $g$ are tangent, meaning their normal vectors (gradients) are parallel.
\begin{formulabox}
    \textbf{Lagrange System of Equations (1 Constraint):} Solve for $x, y, z$ and $\lambda$:
    \[ \begin{cases} \nabla f = \lambda \nabla g \\ g(x,y,z) = k \end{cases} \iff \begin{cases} f_x = \lambda g_x \\ f_y = \lambda g_y \\ f_z = \lambda g_z \\ g(x,y,z) = k \end{cases} \]
\end{formulabox}

\begin{examplebox}{Extrema with Lagrange Multipliers}
    Find the extreme values of $f(x,y) = x^2 + 2y^2$ on the circle $x^2 + y^2 = 1$.\\
    \textit{Solution:}\\
    Let the constraint be $g(x,y) = x^2 + y^2 = 1$. The gradient vectors are:
    \[ \nabla f = \langle 2x, 4y \rangle, \quad \nabla g = \langle 2x, 2y \rangle \]
    Write down the system of equations:
    \begin{align}
        2x &= \lambda (2x) \implies x(1 - \lambda) = 0 \\
        4y &= \lambda (2y) \implies y(2 - \lambda) = 0 \\
        x^2 + y^2 &= 1
    \end{align}
    From equation (1), either $x = 0$ or $\lambda = 1$.
    \begin{itemize}
        \item Case A: $x = 0$. Substituting into (3) yields $y^2 = 1 \implies y = \pm 1$. This gives points $(0, 1)$ and $(0, -1)$.
        \item Case B: $\lambda = 1$. Substituting $\lambda = 1$ into (2) yields $y(2-1) = 0 \implies y = 0$. Substituting into (3) yields $x^2 = 1 \implies x = \pm 1$. This gives points $(1,0)$ and $(-1,0)$.
    \end{itemize}
    Evaluate $f(x,y)$ at these candidate points:
    \begin{itemize}
        \item $f(0, \pm 1) = 0^2 + 2(\pm 1)^2 = 2$ (Absolute Maximum).
        \item $f(\pm 1, 0) = (\pm 1)^2 + 2(0)^2 = 1$ (Absolute Minimum).
    \end{itemize}
\end{examplebox}

\newpage

% =========================================================================
% CHAPTER 2
% =========================================================================
\section{Chapter 2: Multiple Integration}

\subsection{Double Integrals over Rectangles}
\subsubsection*{Key Notations}
\begin{itemize}
    \item $\iint_R f(x,y)\,dA$: Double integral of $f$ over rectangle $R$.
    \item $dA = dy\,dx$ or $dx\,dy$: Area element.
\end{itemize}

\subsubsection*{Concise Explanation \& Formulas}
A double integral over a rectangle sums function values weighted by sub-rectangle areas. Fubini's Theorem states that we can evaluate a double integral over a rectangle as an iterated single integral in either order.
\begin{theorembox}[title={Fubini's Theorem for Rectangles}]
    If $f(x,y)$ is continuous on the rectangular region $R = [a,b] \times [c,d]$, then:
    \[ \iint_R f(x,y)\,dA = \int_a^b \int_c^d f(x,y)\,dy\,dx = \int_c^d \int_a^b f(x,y)\,dx\,dy \]
    If the integrand is separable, i.e., $f(x,y) = g(x)h(y)$, then:
    \[ \iint_R g(x)h(y)\,dA = \left( \int_a^b g(x)\,dx \right) \left( \int_c^d h(y)\,dy \right) \]
\end{theorembox}

\begin{examplebox}{Evaluating Iterated Integrals}
    Evaluate $\iint_R (x - 3y^2)\,dA$ where $R = \{(x,y) \in \mathbb{R}^2 \mid 0 \le x \le 2, \, 1 \le y \le 2\}$.\\
    \textit{Solution:}\\
    Using Fubini's Theorem, integrate first with respect to $y$:
    \begin{align*}
        \iint_R (x - 3y^2)\,dA &= \int_0^2 \int_1^2 (x - 3y^2)\,dy\,dx = \int_0^2 \left[ xy - y^3 \right]_{y=1}^{y=2} dx \\
        &= \int_0^2 \left( (2x - 8) - (x - 1) \right) dx = \int_0^2 (x - 7)\,dx \\
        &= \left[ \frac{1}{2}x^2 - 7x \right]_0^2 = \left(\frac{1}{2}(4) - 14\right) - 0 = 2 - 14 = -12
    \end{align*}
\end{examplebox}


\subsection{Double Integrals over General Regions}
\subsubsection*{Key Notations}
\begin{itemize}
    \item Type I Region: Bounded above and below by continuous functions of $x$.
    \item Type II Region: Bounded left and right by continuous functions of $y$.
\end{itemize}

\subsubsection*{Concise Explanation \& Formulas}
For regions that are not simple rectangles, we set the variable bounds of the inner integral to be functions of the outer variable.
\begin{formulabox}
    \textbf{Type I Region:} $D = \{(x,y) \mid a \le x \le b, \, g_1(x) \le y \le g_2(x)\}$
    \[ \iint_D f(x,y)\,dA = \int_a^b \int_{g_1(x)}^{g_2(x)} f(x,y)\,dy\,dx \]
    \textbf{Type II Region:} $D = \{(x,y) \mid c \le y \le d, \, h_1(y) \le x \le h_2(y)\}$
    \[ \iint_D f(x,y)\,dA = \int_c^d \int_{h_1(y)}^{h_2(y)} f(x,y)\,dx\,dy \]
    \textbf{Area of General Region:}
    \[ A(D) = \iint_D 1\,dA \]
\end{formulabox}

\begin{examplebox}{Double Integral over General Region}
    Evaluate $\iint_D (x + 2y)\,dA$, where $D$ is the region bounded by the parabolas $y = 2x^2$ and $y = 1 + x^2$.\\
    \textit{Solution:}\\
    Step 1: Find the intersection points of the boundaries to get the bounds for $x$.
    \[ 2x^2 = 1 + x^2 \implies x^2 = 1 \implies x = \pm 1 \]
    For $x \in [-1, 1]$, we have $1+x^2 \ge 2x^2$. Thus, $D$ is a Type I region described by:
    \[ D = \{(x,y) \mid -1 \le x \le 1, \, 2x^2 \le y \le 1+x^2\} \]
    Step 2: Set up and evaluate the integral.
    \begin{align*}
        \iint_D (x + 2y)\,dA &= \int_{-1}^1 \int_{2x^2}^{1+x^2} (x + 2y)\,dy\,dx = \int_{-1}^1 \left[ xy + y^2 \right]_{y=2x^2}^{y=1+x^2} dx \\
        &= \int_{-1}^1 \left( x(1+x^2) + (1+x^2)^2 - x(2x^2) - (2x^2)^2 \right) dx \\
        &= \int_{-1}^1 \left( x + x^3 + 1 + 2x^2 + x^4 - 2x^3 - 4x^4 \right) dx \\
        &= \int_{-1}^1 \left( -3x^4 - x^3 + 2x^2 + x + 1 \right) dx
    \end{align*}
    By symmetry, integrals of odd functions ($x^3, x$) over symmetric intervals $[-1, 1]$ are zero. Hence:
    \[ \int_{-1}^1 \left( -3x^4 - x^3 + 2x^2 + x + 1 \right) dx = 2\int_0^1 \left( -3x^4 + 2x^2 + 1 \right) dx \]
    \[ = 2\left[ -\frac{3}{5}x^5 + \frac{2}{3}x^3 + x \right]_0^1 = 2\left( -\frac{3}{5} + \frac{2}{3} + 1 \right) = 2\left( \frac{-9 + 10 + 15}{15} \right) = \frac{32}{15} \]
\end{examplebox}


\subsection{Double Integrals in Polar Coordinates}
\subsubsection*{Key Notations}
\begin{itemize}
    \item $(r, \theta)$: Polar coordinates.
    \item $dA = r\,dr\,d\theta$: The area element in polar coordinates. (Do not forget the extra factor of $r$!)
\end{itemize}

\subsubsection*{Concise Explanation \& Formulas}
Polar coordinates are useful when boundaries of a region are circular or are parameterized naturally by distance from the origin $r$ and angle $\theta$.
\begin{formulabox}
    \textbf{Coordinate Transformations:}
    \[ x = r\cos\theta, \quad y = r\sin\theta, \quad x^2 + y^2 = r^2 \]
    \textbf{Polar Double Integral (General Region):}
    If $D = \{(r,\theta) \mid \alpha \le \theta \le \beta, \, h_1(\theta) \le r \le h_2(\theta)\}$, then:
    \[ \iint_D f(x,y)\,dA = \int_\alpha^\beta \int_{h_1(\theta)}^{h_2(\theta)} f(r\cos\theta, r\sin\theta) \, r\,dr\,d\theta \]
\end{formulabox}

\begin{examplebox}{Double Integral in Polar Coordinates}
    Find the volume of the solid that lies under the paraboloid $z = x^2 + y^2,$, above the $xy$-plane, and inside the cylinder $x^2 + y^2 = 2x$.\\
    \textit{Solution:}\\
    The cylinder is $x^2 + y^2 = 2x$. In polar coordinates, this is:
    \[ r^2 = 2r\cos\theta \implies r = 2\cos\theta \]
    To lie inside the cylinder, $r$ goes from $0$ to $2\cos\theta$. For $r \ge 0$, we require $\cos\theta \ge 0 \implies \theta \in [-\pi/2, \pi/2]$. Thus, the polar region of integration $D$ is:
    \[ D = \{(r,\theta) \mid -\pi/2 \le \theta \le \pi/2, \, 0 \le r \le 2\cos\theta\} \]
    The volume is the integral of the height $z = x^2 + y^2 = r^2$:
    \[ V = \iint_D (x^2+y^2)\,dA = \int
_{-\pi/2}^{\pi/2} \int_0^{2\cos\theta} (r^2) \, r\,dr\,d\theta = \int_{-\pi/2}^{\pi/2} \int_0^{2\cos\theta} r^3 \,dr\,d\theta \]
    \[ = \int_{-\pi/2}^{\pi/2} \left[ \frac{1}{4}r^4 \right]_0^{2\cos\theta} d\theta = \int_{-\pi/2}^{\pi/2} 4\cos^4\theta\,d\theta \]
    Using symmetry and the half-angle identity $\cos^2\theta = \frac{1+\cos(2\theta)}{2}$:
    \[ \cos^4\theta = \left(\frac{1+\cos(2\theta)}{2}\right)^2 = \frac{1}{4}(1 + 2\cos(2\theta) + \cos^2(2\theta)) = \frac{1}{4}\left(1 + 2\cos(2\theta) + \frac{1+\cos(4\theta)}{2}\right) \]
    Therefore:
    \[ 4\cos^4\theta = \frac{3}{2} + 2\cos(2\theta) + \frac{1}{2}\cos(4\theta) \]
    Now integrate:
    \[ V = \int_{-\pi/2}^{\pi/2} \left( \frac{3}{2} + 2\cos(2\theta) + \frac{1}{2}\cos(4\theta) \right) d\theta = \left[ \frac{3}{2}\theta + \sin(2\theta) + \frac{1}{8}\sin(4\theta) \right]_{-\pi/2}^{\pi/2} = \frac{3}{2}\pi \]
\end{examplebox}


\subsection{Triple Integrals}
\subsubsection*{Key Notations}
\begin{itemize}
    \item $\iiint_E f(x,y,z)\,dV$: Triple integral over a 3D solid $E$.
    \item $dV = dz\,dy\,dx$: Volume element in Cartesian coordinates.
    \item $V(E)$: Volume of solid $E$.
\end{itemize}

\subsubsection*{Concise Explanation \& Formulas}
Triple integrals compute quantities (like mass or volume) over 3D domains. For a Type 1 solid region, we integrate along $z$ first, between lower and upper bounding surfaces $u_1(x,y)$ and $u_2(x,y)$, and then integrate over the 2D projection $D$ on the $xy$-plane.
\begin{formulabox}
    \textbf{Fubini's Theorem for a Rectangular Box $B = [a,b] \times [c,d] \times [r,s]$:}
    \[ \iiint_B f(x,y,z)\,dV = \int_r^s \int_c^d \int_a^b f(x,y,z)\,dx\,dy\,dz \]
    \textbf{Integration over Type 1 Solid Region:}
    If $E = \{(x,y,z) \mid (x,y) \in D, \, u_1(x,y) \le z \le u_2(x,y)\}$, then:
    \[ \iiint_E f(x,y,z)\,dV = \iint_D \left( \int_{u_1(x,y)}^{u_2(x,y)} f(x,y,z)\,dz \right) dA \]
    \textbf{Volume of a Solid:}
    \[ V(E) = \iiint_E 1\,dV \]
\end{formulabox}

\begin{examplebox}{Triple Integral and Tetrahedron Volume}
    Evaluate $\iiint_E z\,dV$, where $E$ is the solid tetrahedron bounded by the coordinate planes $x = 0$, $y = 0$, $z = 0$, and the plane $x + y + z = 1$.\\
    \textit{Solution:}\\
    The upper surface is $z = 1 - x - y$ and the lower surface is $z = 0$.
    The projection $D$ on the $xy$-plane is bounded by $x=0$, $y=0$, and $x+y=1 \implies y = 1-x$.
    Thus, the region $E$ is described by:
    \[ E = \{(x,y,z) \mid 0 \le x \le 1, \, 0 \le y \le 1-x, \, 0 \le z \le 1-x-y\} \]
    Set up and evaluate the triple integral:
    \begin{align*}
        \iiint_E z\,dV &= \int_0^1 \int_0^{1-x} \int_0^{1-x-y} z\,dz\,dy\,dx = \int_0^1 \int_0^{1-x} \left[ \frac{1}{2}z^2 \right]_0^{1-x-y} dy\,dx \\
        &= \frac{1}{2} \int_0^1 \int_0^{1-x} (1-x-y)^2\,dy\,dx
    \end{align*}
    Let $u = 1-x-y \implies dy = -du$. When $y=0, u=1-x$. When $y=1-x, u=0$.
    \[ \int_0^{1-x} (1-x-y)^2\,dy = \int_{1-x}^0 u^2(-du) = \int_0^{1-x} u^2\,du = \left[ \frac{1}{3}u^3 \right]_0^{1-x} = \frac{1}{3}(1-x)^3 \]
    Now compute the outer integral:
    \[ \iiint_E z\,dV = \frac{1}{2} \int_0^1 \frac{1}{3}(1-x)^3\,dx = \frac{1}{6} \left[ -\frac{1}{4}(1-x)^4 \right]_0^1 = \frac{1}{6}\left( 0 - \left(-\frac{1}{4}\right) \right) = \frac{1}{24} \]
\end{examplebox}


\subsection{Triple Integrals in Cylindrical Coordinates}
\subsubsection*{Key Notations}
\begin{itemize}
    \item $(r, \theta, z)$: Cylindrical coordinates.
    \item $dV = r\,dz\,dr\,d\theta$: The volume element. (Includes the scaling factor $r$).
\end{itemize}

\subsubsection*{Concise Explanation \& Formulas}
Cylindrical coordinates combine polar coordinates in the $xy$-plane with standard linear height $z$. They are ideal for solids with axial symmetry around the $z$-axis.
\begin{formulabox}
    \textbf{Coordinate Transformations:}
    \[ x = r\cos\theta, \quad y = r\sin\theta, \quad z = z, \quad x^2 + y^2 = r^2 \]
    \textbf{Cylindrical Integration Formula:}
    \[ \iiint_E f(x,y,z)\,dV = \int_\alpha^\beta \int_{h_1(\theta)}^{h_2(\theta)} \int_{u_1(r\cos\theta, r\sin\theta)}^{u_2(r\cos\theta, r\sin\theta)} f(r\cos\theta, r\sin\theta, z) \, r\,dz\,dr\,d\theta \]
\end{formulabox}

\begin{examplebox}{Triple Integral in Cylindrical Coordinates}
    Evaluate $\iiint_E \sqrt{x^2+y^2}\,dV$, where $E$ is the region that lies within the cylinder $x^2 + y^2 = 1$, below the plane $z = 4$, and above the paraboloid $z = 1 - x^2 - y^2$.\\
    \textit{Solution:}\\
    In cylindrical coordinates, the boundaries are $r=1$ (cylinder), $z=4$ (top plane), and $z=1-r^2$ (paraboloid).
    The solid region is $E = \{(r,\theta,z) \mid 0 \le \theta \le 2\pi, \, 0 \le r \le 1, \, 1-r^2 \le z \le 4\}$.
    The integrand is $\sqrt{x^2+y^2} = r$.
    Set up the integral:
    \begin{align*}
        \iiint_E \sqrt{x^2+y^2}\,dV &= \int_0^{2\pi} \int_0^1 \int_{1-r^2}^4 (r) \, r\,dz\,dr\,d\theta = \int_0^{2\pi} \int_0^1 r^2 \left( 4 - (1-r^2) \right) dr\,d\theta \\
        &= 2\pi \int_0^1 r^2(3+r^2)\,dr = 2\pi \int_0^1 (3r^2 + r^4)\,dr \\
        &= 2\pi \left[ r^3 + \frac{1}{5}r^5 \right]_0^1 = 2\pi\left(1 + \frac{1}{5}\right) = \frac{12\pi}{5}
    \end{align*}
\end{examplebox}


\subsection{Triple Integrals in Spherical Coordinates}
\subsubsection*{Key Notations}
\begin{itemize}
    \item $(\rho, \theta, \phi)$: Spherical coordinates.
    \item $\rho$: Distance from origin ($\rho \ge 0$).
    \item $\theta$: Angle in the $xy$-plane ($0 \le \theta \le 2\pi$).
    \item $\phi$: Angle from the positive $z$-axis ($0 \le \phi \le \pi$).
    \item $dV = \rho^2\sin\phi\,d\rho\,d\theta\,d\phi$: The volume element. (Includes scaling factor $\rho^2\sin\phi$).
\end{itemize}

\subsubsection*{Concise Explanation \& Formulas}
Spherical coordinates define a point by distance $\rho$ and two angles $\theta$ and $\phi$. They are highly efficient for spheres, cones, or regions symmetrical about the origin.
\begin{formulabox}
    \textbf{Coordinate Transformations:}
    \[ x = \rho\sin\phi\cos\theta, \quad y = \rho\sin\phi\sin\theta, \quad z = \rho\cos\phi, \quad x^2 + y^2 + z^2 = \rho^2 \]
    \textbf{Spherical Integration Formula (over a wedge):}
    If $E = \{(\rho, \theta, \phi) \mid a \le \rho \le b, \, \alpha \le \theta \le \beta, \, c \le \phi \le d\}$, then:
    \[ \iiint_E f(x,y,z)\,dV = \int_c^d \int_\alpha^\beta \int_a^b f(\rho\sin\phi\cos\theta, \dots) \, \rho^2\sin\phi\,d\rho\,d\theta\,d\phi \]
\end{formulabox}

\begin{examplebox}{Volume using Spherical Coordinates}
    Find the volume of the solid $E$ that lies above the cone $z = \sqrt{x^2+y^2}$ and below the sphere $x^2+y^2+z^2 = z$.\\
    \textit{Solution:}\\
    Step 1: Write the boundaries in spherical coordinates.
    \begin{itemize}
        \item Sphere: $\rho^2 = \rho\cos\phi \implies \rho = \cos\phi$.
        \item Cone: $\rho\cos\phi = \sqrt{\rho^2\sin^2\phi(\cos^2\theta+\sin^2\theta)} = \rho\sin\phi \implies \tan\phi = 1 \implies \phi = \pi/4$.
    \end{itemize}
    Since the solid lies above the cone, the angle $\phi$ ranges from $0$ (the positive $z$-axis) to $\pi/4$.
    For any angle $\phi$, $\rho$ ranges from $0$ (origin) to $\cos\phi$.
    The angle $\theta$ ranges from $0$ to $2\pi$ around the $z$-axis.
    Thus, $E = \{(\rho,\theta,\phi) \mid 0 \le \theta \le 2\pi, \, 0 \le \phi \le \pi/4, \, 0 \le \rho \le \cos\phi\}$.
    Step 2: Integrate to find volume.
    \begin{align*}
        V(E) &= \iiint_E dV = \int_0^{2\pi} \int_0^{\pi/4} \int_0^{\cos\phi} \rho^2\sin\phi\,d\rho\,d\phi\,d\theta \\
        &= 2\pi \int_0^{\pi/4} \sin\phi \left[ \frac{1}{3}\rho^3 \right]_0^{\cos\phi} d\phi = \frac{2\pi}{3} \int_0^{\pi/4} \cos^3\phi\sin\phi\,d\phi
    \end{align*}
    Let $u = \cos\phi \implies du = -\sin\phi\,d\phi$.
    When $\phi = 0, u = 1$. When $\phi = \pi/4, u = 1/\sqrt{2}$.
    \[ V(E) = \frac{2\pi}{3} \int_{1/\sqrt{2}}^1 u^3\,du = \frac{2\pi}{3} \left[ \frac{1}{4}u^4 \right]_{1/2}^1 = \frac{\pi}{6} \left( 1 - \frac{1}{4} \right) = \frac{\pi}{6}\left(\frac{3}{4}\right) = \frac{\pi}{8} \]
\end{examplebox}

\end{document}
